\subsection{Homomorphisms compatible with filtrations}

Let \(G, G'\) be two commutative groups (written additively), \((G_n)\) a filtration on \(G\) and \((G'_n)\) a filtration on \(G'\) ; a homomorphism \(h: G \to G'\) is called compatible with the filtrations on \(G\) and \(G'\) if \(h(G_n) \subset G'_n\) for all \(n \in \mathbf{Z}\) . The composite homomorphism \(G_n \xrightarrow{h|G_n} G'_n \longrightarrow G'_n / G_{n+1}'\) is zero on \(G_{n+1}\) and hence defines by taking quotients a homomorphism \(h_n: G_n / G_{n+1} \to G'_n / G_{n+1}'\) ; there is therefore a unique additive group homomorphism \(\text{gr}(h): \text{gr}(G) \to \text{gr}(G')\) such that, for all \(n \in \mathbf{Z}\) , \(\text{gr}(h)\) coincides with \(h_n\) on \(\text{gr}_n(G) = G_n / G_{n+1}\) . \(\text{gr}(h)\) is called the graded group homomorphism associated with \(h\) . If \(G''\) is a third filtered group and \(h': G' \to G''\) a homomorphism which is compatible with the filtrations, \(h' \circ h\) is a homomorphism which is compatible with the filtrations and

\[
\operatorname{gr} (h ^ {\prime} \circ h) = \operatorname{gr} (h ^ {\prime}) \circ \operatorname{gr} (h).\tag{12}
\]

PROPOSITION 2. Let G be a filtered commutative group and H a subgroup of G; let H be given the induced filtration and G/H the quotient filtration. If j: H → G is the canonical injection and p: G → G/H the canonical surjection, j and p are compatible with the filtrations and the sequence

\[
0 \longrightarrow \operatorname{gr} (\mathrm{H}) \xrightarrow {\operatorname{gr} (j)} \operatorname{gr} (\mathrm{G}) \xrightarrow {\operatorname{gr} (p)} \operatorname{gr} (\mathrm{G} / \mathrm{H}) \longrightarrow 0\tag{13}
\]

is exact.

The first assertion is obvious; if \((\mathbf{G}_n)\) is the filtration on \(\mathbf{G}\) , then

\[
\left(\mathrm{H} \cap \mathrm{G} _ {n}\right) \cap \mathrm{G} _ {n + 1} = \mathrm{H} \cap \mathrm{G} _ {n + 1}
\]

and hence \(\mathrm{gr}(j)\) is injective; moreover the canonical mapping

\[
\mathrm{G} _ {n} \rightarrow (\mathrm{H} + \mathrm{G} _ {n}) / \mathrm{H}
\]

is surjective, hence so is \(\operatorname{gr}(p)\) and \(\operatorname{gr}(p) \circ \operatorname{gr}(j) = 0\) by (12). Finally, let \(\xi \in \mathbf{G}_n / \mathbf{G}_{n+1}\) belong to the kernel of \(\operatorname{gr}(p)\) ; then there exists \(x \in \xi\) such that \(x \in \mathbf{H} + \mathbf{G}_{n+1}\) ; but as \(\mathbf{G}_{n+1} \subset \mathbf{G}_n\) ,

\[
\mathrm{G} _ {n} \cap (\mathrm{H} + \mathrm{G} _ {n + 1}) = (\mathrm{H} \cap \mathrm{G} _ {n}) + \mathrm{G} _ {n + 1}
\]

and hence \(x = y + z\) where \(y \in \mathbf{H} \cap \mathbf{G}_n\) and \(z \in \mathbf{G}_{n+1}\) ; this proves that \(\xi\) is the class mod. \(\mathbf{G}_{n+1}\) of \(j(y)\) , in other words it belongs to the image of \(\operatorname{gr}(\mathbf{H})\) under \(\operatorname{gr}(j)\) .

Note that, given an exact sequence \(0 \to G' \xrightarrow{u} G \xrightarrow{v} G'' \to 0\) of filtered commutative groups, where \(u\) and \(v\) are compatible with the filtrations, the sequence \(0 \longrightarrow \text{gr}(G') \xrightarrow{\text{gr}(u)} \text{gr}(G) \xrightarrow{\text{gr}(v)} \text{gr}(G'') \longrightarrow 0\) is not necessarily exact (Exercise 4).

If now A and B are two filtered rings and \(h \colon \mathbf{A} \to \mathbf{B}\) a ring homomorphism which is compatible with the filtrations, it is immediately verified that the graded group homomorphism \(\mathrm{gr}(h) \colon \mathrm{gr}(\mathbf{A}) \to \mathrm{gr}(\mathbf{B})\) is also a ring homomorphism. In particular, if \(\mathbf{A}'\) is a subring of A with the induced filtration, \(\mathrm{gr}(\mathbf{A}')\) is canonically identified with a graded subring of \(\mathrm{gr}(\mathbf{A})\) (Proposition 2); if \(\mathfrak{b}\) is a two-sided ideal of A and \(\mathbf{A} / \mathfrak{b}\) is given the quotient filtration, \(\mathrm{gr}(\mathbf{A} / \mathfrak{b})\) is canonically identified with the quotient graded ring \(\mathrm{gr}(\mathbf{A}) / \mathrm{gr}(\mathfrak{b})\) (Proposition 2).

Finally, let A be a filtered ring, E, F two filtered A-modules and \(u: \mathbf{E} \to \mathbf{F}\) an homomorphism compatible with the filtrations. Then it is immediate that \(\mathrm{gr}(u): \mathrm{gr}(\mathrm{E}) \to \mathrm{gr}(\mathrm{F})\) is a \(\mathrm{gr}(\mathrm{A})\) -linear mapping and hence a homogeneous homomorphism of degree 0 of graded \(\mathrm{gr}(\mathrm{A})\) -modules. Moreover, if \(u': \mathbf{E} \to \mathbf{F}\) is another A-homomorphism compatible with the filtrations, so is \(u + u'\) and

\[
\operatorname{gr} (u + u ^ {\prime}) = \operatorname{gr} (u) + \operatorname{gr} (u ^ {\prime}).\tag{14}
\]

\textbf{Remarks}\par

\begin{enumerate}
\def\labelenumi{(\arabic{enumi})}
\item
  Clearly filtered ring homomorphisms (resp. homomorphisms of filtered modules over a given filtered ring A) compatible with the filtrations can be taken as morphisms for the filtered ring structure (resp. filtered A-module structure) (Set Theory, Chapter IV, § 2, no. 1).
\item
  Let E and F be two modules over a filtered ring A and let them have the filtrations derived from the filtration \((\mathbf{A}_{n})\) on A (no. 1, Example 2). Then every A-linear mapping \(u: E \rightarrow F\) is compatible with the filtrations, since

\[
u (\mathrm{A} _ {n} \mathrm{E}) = \mathrm{A} _ {n} u (\mathrm{E}) \subset \mathrm{A} _ {n} \mathrm{F}.
\]

\item
  Note that a filtered A-module homomorphism \(u: \mathbf{E} \to \mathbf{F}\) which is compatible with the filtrations may satisfy \(\operatorname{gr}(u) = 0\) without being zero; this is so for example of the endomorphism \(x \mapsto nx\) of the additive group \(\mathbf{Z}\) with the \((n)\) -adic filtration (for any integer \(n > 1\) ). The relation \(\operatorname{gr}(u) = \operatorname{gr}(v)\) for two homomorphisms \(u, v\) of \(\mathbf{E}\) to \(\mathbf{F}\) , compatible with the filtrations, does not therefore imply necessarily \(u = v\) .
\item
  The definitions at the beginning of this no. extend immediately to two groups G, G', which are not necessarily commutative and are filtered by subgroups Gn, Gn' such that Gn+1 (resp. Gn'+1) is normal in Gn (resp. Gn'). Proposition 2 is also valid with the same hypotheses on the Gn and assuming that H is invariant in G, the proof remaining unchanged except for notation.
\end{enumerate}

